% !TeX program = xelatex
% Generated by build_speech_notes_cn.py from speech_notes_cn.json.
\documentclass[11pt,a4paper]{article}
\usepackage[left=18mm,right=18mm,top=17mm,bottom=16mm,headheight=14pt,headsep=5mm]{geometry}
\usepackage{fontspec,xeCJK,graphicx,xcolor,fancyhdr,hyperref}
\setmainfont{texgyretermes}[Extension=.otf,UprightFont=*-regular,BoldFont=*-bold,ItalicFont=*-italic,BoldItalicFont=*-bolditalic]
\setsansfont{texgyreheros}[Extension=.otf,UprightFont=*-regular,BoldFont=*-bold,ItalicFont=*-italic,BoldItalicFont=*-bolditalic]
\setCJKmainfont{Noto Serif CJK SC}
\setCJKsansfont{Noto Sans CJK SC}
\definecolor{Navy}{HTML}{17324D}
\definecolor{Teal}{HTML}{007F82}
\definecolor{Slate}{HTML}{586878}
\pagestyle{fancy}\fancyhf{}
\fancyhead[L]{\sffamily\small\color{Teal}TRACE-Q 中文逐页讲稿}
\fancyhead[R]{\small\color{Slate}对应原版 17 页幻灯片}
\fancyfoot[L]{\fontsize{8}{10}\selectfont\color{Slate}正文与转场可直接口述；讲解提示不必读出。}
\fancyfoot[R]{\small\color{Slate}\thepage\,/\,17}
\renewcommand{\headrulewidth}{0.3pt}
\renewcommand{\footrulewidth}{0pt}
\setlength{\parindent}{0pt}
\setlength{\parskip}{5pt}
\setlength{\emergencystretch}{2em}
\hypersetup{unicode=true,colorlinks=true,linkcolor=Teal,
pdftitle={TRACE-Q 中文逐页讲稿},pdfsubject={17 页幻灯片的中文口述稿与转场，含原幻灯片缩略图}}
\begin{document}
\fontsize{11}{16}\selectfont

\pdfbookmark[1]{第 1 页：汇报路线：研究设计、审查规则与时期变化}{slide-1}
{\sffamily\bfseries\color{Navy}\fontsize{14}{19}\selectfont 第 1 页｜汇报路线：研究设计、审查规则与时期变化\par}
{\sffamily\color{Slate}\fontsize{8.5}{11}\selectfont TRACE-Q: study design, review evidence and period changes\par}
\IfFileExists{trace_q_slides.pdf}{\begin{center}\includegraphics[page=1,width=.76\linewidth]{trace_q_slides.pdf}\end{center}}{}
{\color{Slate}\fontsize{9}{12}\selectfont\textbf{讲解提示：}先说明三个部分；全部确认与确认且实质是包含关系，不能相加。\par}
\vspace{1mm}
今天的汇报以 review\_results\_2\_4\_1\_Rbulun\_copy 文件夹中四份 completed 工作簿的最终审查结果为依据。第一部分介绍 TRACE-Q 研究设计与论文样本；第二部分说明 codebook 如何把候选条目转化为可汇总的判断；第三部分比较早期与后期发表队列。原始说明的机制分类作为补充分析单独介绍。\par
我们重新读取了四份工作簿的候选工作表，逐项核对处置、最终领域、完成状态、证据判断和实质性。统计结果中，已确认的问题共 139 个，其中 37 个具有实质性。37 是 139 的子集，不能相加。下面的数量和时期比较都以这些最终判断为准。\par
这里的疫情前后，具体指 2018 至 2019 年和 2024 至 2025 年发表的论文，不是按每个试验的招募时间分组，也不是疫情影响的因果研究。文本分析也有明确边界：只重新分类 statistical\_results 中已确认的问题。其他三个领域仅用于单独的时期比较。\par
{\color{Teal}\textbf{转场：}先介绍实际纳入的论文和审查资料。\par}
\vfill{\color{Slate}\fontsize{8.5}{11}\selectfont 可按听众背景调整研究设计、codebook 与结果部分的讲解时间。\par}
\clearpage
\pdfbookmark[1]{第 2 页：样本：44 个试验资料包与最终人工处置}{slide-2}
{\sffamily\bfseries\color{Navy}\fontsize{14}{19}\selectfont 第 2 页｜样本：44 个试验资料包与最终人工处置\par}
{\sffamily\color{Slate}\fontsize{8.5}{11}\selectfont The comparison uses 44 selected trial packages\par}
\IfFileExists{trace_q_slides.pdf}{\begin{center}\includegraphics[page=2,width=.76\linewidth]{trace_q_slides.pdf}\end{center}}{}
{\color{Slate}\fontsize{9}{12}\selectfont\textbf{讲解提示：}重点报 19 对 25、440 对 413；表内各年的页数无需逐项朗读。\par}
\vspace{1mm}
四份最终工作簿对应 44 篇论文：2018 年 10 篇，2019 年 9 篇，2024 年 10 篇，2025 年 15 篇。早期是 19 篇，后期是 25 篇。年份按工作簿的候选编号和 Papers 清单确定，不能只根据 DOI 中的年份分组。\par
逐年候选数是 115、79、113 和 133；保留数是 107、76、107 和 123。合起来有 440 个候选、413 个保留和 27 个排除。每个候选的完成状态都是 primary\_done，没有 discuss 处置。保留并不等于已确认统计错误，各领域还要按自己的判断字段继续筛选。\par
这页的 188 个 PDF、5,559 页和补充 Excel 是原有资料清单描述的材料规模，并不是本次从审查工作簿重新计算的结果。核心分析的分母仍然是 44 篇论文，不能把每一页或每个候选当成独立的论文观察值。\par
研究采用资料可获得性样本，不能代表期刊所有随机试验。工作簿提供最终判断，但现有资料没有可识别的独立人工双人评分，因此不能用这些结果估计评审者间一致性。\par
还有一个影响讲稿解释的事实：最终工作簿中，413 条保留记录的 Notes 全部为空；非空说明只出现在 27 条排除记录中。因此，后面的确认数和实质性数可以直接核对，机制类别和具体解释则不能声称来自这些最终工作簿的说明栏。\par
{\color{Teal}\textbf{转场：}接着用 codebook 说明，保留一个条目究竟代表什么。\par}
\clearpage
\pdfbookmark[1]{第 3 页：Codebook：先确定处置和分支，再完成判断}{slide-3}
{\sffamily\bfseries\color{Navy}\fontsize{14}{19}\selectfont 第 3 页｜Codebook：先确定处置和分支，再完成判断\par}
{\sffamily\color{Slate}\fontsize{8.5}{11}\selectfont Common fields establish a traceable human decision\par}
\IfFileExists{trace_q_slides.pdf}{\begin{center}\includegraphics[page=3,width=.76\linewidth]{trace_q_slides.pdf}\end{center}}{}
{\color{Slate}\fontsize{9}{12}\selectfont\textbf{讲解提示：}按“身份与证据—分支—处置—完成状态”解释，不必逐个读 Excel 列号。\par}
\vspace{1mm}
这一页介绍所有候选条目共用的字段。首先，每个条目都有固定的身份编号、标题和来源位置。编号用于持续追踪，不能随意修改；标题只是一条待审查的线索，判断时仍然要打开它指向的原始文件和页码。\par
然后选择审查分支。AI 提供的领域标签可以由人工修改，但原始文件之间的关系代码保持不变。这样既保留人的判断，也不会丢失这个疑点最初来自哪些材料。\par
处置有三个选项：retain 是保留，exclude 是排除，discuss 是已经考虑过、但需要进一步讨论。空白则表示没有回答，不能把空白当成待讨论，也不能当成没有问题。\par
完成状态是另一个维度。起草时用 no，人工检查完成后才标记 primary\_done。如果要改动已经完成的条目，应先恢复到未完成，修改后重新核对。PRIMARY READY 只是表单有效、人工录入完成的提示；即使处置是待讨论，也可能显示这个提示，它不是科学正确性的认证。\par
最后，不同分支要求填写的内容不同。保留的统计条目必须同时判断问题是否成立和是否具有实质性；报告内部或计划材料条目必须判断不一致状态；方案与报告之间的差异则允许只保留描述，解释字段是可选的。\par
{\color{Teal}\textbf{转场：}我们先看统计条目的第一个必填判断：现有证据究竟能确认到什么程度。\par}
\clearpage
\pdfbookmark[1]{第 4 页：证据确定程度：确认、很可能、不是问题或无法判断}{slide-4}
{\sffamily\bfseries\color{Navy}\fontsize{14}{19}\selectfont 第 4 页｜证据确定程度：确认、很可能、不是问题或无法判断\par}
{\sffamily\color{Slate}\fontsize{8.5}{11}\selectfont Statistical evidence: certainty is its own judgment\par}
\IfFileExists{trace_q_slides.pdf}{\begin{center}\includegraphics[page=4,width=.76\linewidth]{trace_q_slides.pdf}\end{center}}{}
{\color{Slate}\fontsize{9}{12}\selectfont\textbf{讲解提示：}先解释四种状态，再讲匹配分析背景；读到 5\% 时明确它不是 P 值阈值。\par}
\vspace{1mm}
统计条目的第一项判断，是是否存在统计或数值报告问题。这里有四种有效状态。confirmed\_issue 表示现有证据已经足以建立一个报告问题或矛盾。一个重要细节是：即使不知道正确的替代数值，也可能确认原来的报告存在矛盾。\par
probable\_issue 表示证据倾向于有问题，但还存在合理的替代解释或未解决的限制。not\_issue 表示核实计算方法、定义或分析背景以后，原来的疑点可以得到解释。indeterminate 则表示证据不足，既不能确认，也不能排除。\par
作出判断之前，必须先确认比较的是同一人群、结局、时间点、比较对象和模型。例如，不能直接拿原始事件数算出的粗比值，去否定一个经过模型调整的估计值。印刷时的四舍五入，以及来源明确说明的不同分析方法，也都要考虑。\par
幻灯片下方还有一条工作流的数值筛选规则：先根据显示精度给输入值留出半个最小显示单位的舍入范围，再传播这个范围。剩余差距要超过比较值绝对值的 5\%，才通过这项数值筛选；比较值为零时，任何正的剩余差距都通过。区间与检验或 P 值不相容的检查不受这项限制。这个 5\% 是操作性筛选规则，与显著性水平 0.05 是两回事。\par
{\color{Teal}\textbf{转场：}确定存在问题以后，还要继续问：这个问题到底只是呈现上的瑕疵，还是可能影响结果的使用。\par}
\clearpage
\pdfbookmark[1]{第 5 页：实质性与下游影响：分别回答两个问题}{slide-5}
{\sffamily\bfseries\color{Navy}\fontsize{14}{19}\selectfont 第 5 页｜实质性与下游影响：分别回答两个问题\par}
{\sffamily\color{Slate}\fontsize{8.5}{11}\selectfont Substantiveness and potential impact are separate\par}
\IfFileExists{trace_q_slides.pdf}{\begin{center}\includegraphics[page=5,width=.76\linewidth]{trace_q_slides.pdf}\end{center}}{}
{\color{Slate}\fontsize{9}{12}\selectfont\textbf{讲解提示：}先指左侧必填项，再指右侧可选项；强调 direct 表示潜在使用路径。\par}
\vspace{1mm}
这一页把两个容易混淆的概念分开。左边的实质性是必填项。substantive 表示这个问题可能实质性地影响结果解释、数据提取或者后续分析。minor\_presentation 表示主要是小的格式、标签或印刷问题，读者仍然能清楚理解原本的结果。indeterminate 表示还无法区分这两类。\par
比如，一个无法确定该提取哪个数值的效应量，可能属于实质性问题；一个不影响理解的空格错误，则属于呈现问题。所以，“确认”说的是证据是否足够，“实质性”说的是这个问题可能有多大影响，两者不能互相替代。\par
右边的下游影响是可选项。direct 表示受影响的数值或定义本身就是后续提取或分析的输入；possible 表示存在可能的影响路径，但比较间接或不确定；unlikely 表示相关结果仍然清楚可用，不太可能影响后续使用；unknown 表示现有材料不足以评估。\par
这里的 direct 并不表示已经观察到后续分析受到损害。可选字段留空表示没有评估，not\_applicable 则表示这个字段在当前分支不适用。这项可选判断不决定是否进入后面的时期比较。\par
{\color{Teal}\textbf{转场：}统计条目需要这两个独立判断，其他类型的条目则采用不同的评价方式。\par}
\clearpage
\pdfbookmark[1]{第 6 页：其他领域：报告内部、计划材料与方案—报告差异}{slide-6}
{\sffamily\bfseries\color{Navy}\fontsize{14}{19}\selectfont 第 6 页｜其他领域：报告内部、计划材料与方案—报告差异\par}
{\sffamily\color{Slate}\fontsize{8.5}{11}\selectfont Other domains use different assessment branches\par}
\IfFileExists{trace_q_slides.pdf}{\begin{center}\includegraphics[page=6,width=.76\linewidth]{trace_q_slides.pdf}\end{center}}{}
{\color{Slate}\fontsize{9}{12}\selectfont\textbf{讲解提示：}左右分支分别讲解；文件关系代码用于说明材料来源，不必反复逐字读。\par}
\vspace{1mm}
左侧包括报告内部不一致和计划材料不一致。前者比较正文或补充结果材料中的内容，后者比较研究方案、统计分析计划等计划性材料。这里要求判断的是：存在实质性不一致、仅有呈现不一致、实际上没有不一致，还是无法判断。\par
实质性不一致，是指相互冲突的描述可能改变我们对研究、人群或结果的理解。呈现不一致则是真有差别，但本意仍然清楚。如果核对定义、人群和时间点以后发现两种说法相容，就应判断为没有不一致。材料不足时使用无法判断。\par
右侧是方案与报告之间的比较。它的解释字段是可选的，可以记录实质性差异、轻微呈现差异、版本或时间顺序不明、没有不一致，或者无法判断。比如，计划和报告采用了不同的结局定义，可以先记录这个差异；仅凭这一比较，还不能断言修改没有正当理由。当前 codebook 也不要求为了完成这个字段，额外搜索所有修改记录和原因。\par
文件关系代码说明证据来源：S\_C 是报告文件之间，S\_W 是同一报告内部，P\_P 是计划材料内部或之间，P\_R 是计划与报告之间。论文原有的报告材料结局继续排除原始 P\_R 和 P\_P。此次新增的另外三个领域比较，按最终领域标签单独定义，因此允许纳入这些关系；它们是明确标注的描述性扩展。\par
{\color{Teal}\textbf{转场：}这些领域的状态不同，因此后面的比较必须逐一明确计数口径。\par}
\clearpage
\pdfbookmark[1]{第 7 页：分析口径：139、37、41 分别是怎么得到的}{slide-7}
{\sffamily\bfseries\color{Navy}\fontsize{14}{19}\selectfont 第 7 页｜分析口径：139、37、41 分别是怎么得到的\par}
{\sffamily\color{Slate}\fontsize{8.5}{11}\selectfont Translate the codebook into explicit endpoints\par}
\IfFileExists{trace_q_slides.pdf}{\begin{center}\includegraphics[page=7,width=.76\linewidth]{trace_q_slides.pdf}\end{center}}{}
{\color{Slate}\fontsize{9}{12}\selectfont\textbf{讲解提示：}按共享筛选条件、三行结局、149 的区别依次讲；这是连接 codebook 与结果的关键页。\par}
\vspace{1mm}
这一页把最终工作簿与分析口径连接起来。我们读取 D 列处置、E 列最终领域、H 列统计问题判断和 I 列实质性，先筛选 retain 和 statistical\_results。论文原有口径还要求原始文件关系为 S\_C 或 S\_W；关系代码来自保留的来源对照表，最终工作簿没有单独的关系代码列。\par
再要求 H 列是 confirmed\_issue，得到 139 个已确认统计问题。进一步要求 I 列是 substantive，得到 37 个；其余 102 个确认条目的实质性标签为 minor\_presentation。逐年全部确认数为 34、22、43、40；其中实质性为 11、5、11、10。\par
如果保留实质性条件，同时把 probable\_issue 纳入，就得到论文主要结局的 41 个，即 37 个确认加 4 个很可能存在的问题。这四个 probable 都在早期队列。三个口径使用相同的最终工作簿，只是明确选择了不同的证据纳入条件。\par
另外，补充材料里还有 149 个已确认的报告材料问题。这个数更宽，等于 139 个统计问题，再加 10 个符合条件的报告内部不一致问题。它不能替代这里的 139。无论采用哪个口径，完成标记和下游影响评分，都不能代替证据确定程度或实质性判断。\par
{\color{Teal}\textbf{转场：}下面解释原始说明如何与这些最终判断连接。\par}
\clearpage
\pdfbookmark[1]{第 8 页：结果依据与补充说明：最终工作簿不含保留条目的 Notes}{slide-8}
{\sffamily\bfseries\color{Navy}\fontsize{14}{19}\selectfont 第 8 页｜结果依据与补充说明：最终工作簿不含保留条目的 Notes\par}
{\sffamily\color{Slate}\fontsize{8.5}{11}\selectfont Recover raw notes, then add a separate mechanism taxonomy\par}
\IfFileExists{trace_q_slides.pdf}{\begin{center}\includegraphics[page=8,width=.76\linewidth]{trace_q_slides.pdf}\end{center}}{}
{\color{Slate}\fontsize{9}{12}\selectfont\textbf{讲解提示：}明确区分最终判断与原始说明；440 条非空说明属于先前归档来源。\par}
\vspace{1mm}
先区分两种资料：指定文件夹里的最终 completed 工作簿，是处置、领域、证据状态和实质性的依据；先前原始工作簿中的 Notes，是另一个来源。当前最终工作簿共有 440 条候选，但只有 27 条排除记录有说明，413 条保留记录的说明为空。\par
因此，幻灯片中的 440 条非空说明，是先前从原始 R\_Bulun 工作簿归档的提取结果，不能描述为从本次指定文件夹提取的 440 条人工评语。这些原始说明由较早的 AI 审查生成，保留在补充资料中用于解释机制。\par
本次用候选编号重新连接最终判断，确认补充分类中的 139 条全部属于最终保留且确认的统计问题，37 条实质性标签也逐项一致。最终标签决定哪些条目可以进入分析，原始说明不覆盖或补填最终判断。\par
后两页涉及的七类机制，是对已归档说明所做的额外内容分析。词频、完整语境阅读和第二次模型审计都属于这一补充流程；这些机制不是最终工作簿的原始字段，也不能把说明缺失解释为没有问题。\par
{\color{Teal}\textbf{转场：}接下来展示重新分类的七类机制及其含义。\par}
\clearpage
\pdfbookmark[1]{第 9 页：七类统计机制：不能只看标题或关键词}{slide-9}
{\sffamily\bfseries\color{Navy}\fontsize{14}{19}\selectfont 第 9 页｜七类统计机制：不能只看标题或关键词\par}
{\sffamily\color{Slate}\fontsize{8.5}{11}\selectfont Recategorize confirmed statistical issues by mechanism\par}
\IfFileExists{trace_q_slides.pdf}{\begin{center}\includegraphics[page=9,width=.76\linewidth]{trace_q_slides.pdf}\end{center}}{}
{\color{Slate}\fontsize{9}{12}\selectfont\textbf{讲解提示：}七类机制是补充分类；139 个纳入与 37 个实质性标签已和最终工作簿逐项核对。\par}
\vspace{1mm}
这一页继续展示补充机制分类，不能把七个类别当成最终工作簿直接给出的审查结果。类别依次是人数与分母、效应量与估计目标、区间与分布摘要、P 值与推断、比较方向、结果归属，以及分析或测量定义。每条只计一个主要类别，合计 139 条，其中实质性 37 条由最终工作簿决定。\par
人数与分母类最多，有 39 条，但只有 4 条具有实质性。实质性子集中，效应量类有 8 条，区间类 7 条，归属类 6 条；三类合计 21 条，占 37 条的 56.8\%。类别描述问题机制，实质性描述对解释或使用的影响，两者不能互相替代。\par
右侧两个解释性例子来自先前归档的原始说明。一个指出风险比可由转置后的二乘二表复现，另一个指出区间上限印成 42、原意为 4.2。这些细节不能从当前空白的保留记录 Notes 中恢复，因此在这里作为历史说明的例子，不当作本次重新核验原文后得到的结论。\par
归档分类记录包含第二个模型对 139 条的审计、三个主类别修订以及三个歧义标记。本次核对的是这些候选与最终判断能否一致连接，没有重新进行源论文裁决。机制分类仍是 AI 辅助的补充结果，需要与最终审查标签分开理解。\par
{\color{Teal}\textbf{转场：}理解计数对象后，再看论文层面的时期比较方法。\par}
\clearpage
\pdfbookmark[1]{第 10 页：统计方法：以论文为单位进行置换与重抽样}{slide-10}
{\sffamily\bfseries\color{Navy}\fontsize{14}{19}\selectfont 第 10 页｜统计方法：以论文为单位进行置换与重抽样\par}
{\sffamily\color{Slate}\fontsize{8.5}{11}\selectfont Compare publication cohorts at the paper level\par}
\IfFileExists{trace_q_slides.pdf}{\begin{center}\includegraphics[page=10,width=.76\linewidth]{trace_q_slides.pdf}\end{center}}{}
{\color{Slate}\fontsize{9}{12}\selectfont\textbf{讲解提示：}用“后期减早期”口头解释公式；区分精确置换 P 值与 bootstrap 区间。\par}
\vspace{1mm}
这里的分析单位是论文。同一篇论文中的多个问题共享资料和审查背景，不能把它们当成互相独立的样本。因此，我们给每篇论文计算一个问题数，44 篇全部进入分析，没有符合条件问题的论文也以零计入。\par
左边公式表示后期平均数减去早期平均数：后期是 25 篇，早期是 19 篇。差值大于零，表示后期每篇论文的问题更多；小于零，表示后期更少。我们还分析另一个结局：每篇论文是否至少有一个问题，对应的是受影响论文的比例。\par
P 值来自精确的双侧论文层面置换检验。在零假设下具备可交换性的前提下，把整篇论文重新分到 19 篇和 25 篇的两组中，比较哪些分组产生的差异绝对值至少与观察结果一样大。程序用动态规划计数全部分组，并把与观察值一样极端的情况计入。\par
置信区间则用另一种方法计算：在每个时期内部，以论文为单位进行两万次有放回重抽样，取百分位数区间。它不是置换区间。所有 P 值都是探索性的、未经多重比较校正；论文之间的可交换性也可能受到试验类型和资料复杂度差异的影响。\par
本次直接读取最终工作簿后，所有论文层面的计数均与原分析一致，并重新核对了 32 个精确置换检验。置信区间沿用相同论文计数下的既有 20,000 次 bootstrap 结果，没有把它说成一次新的重抽样分析。\par
{\color{Teal}\textbf{转场：}先看两个核心统计结局的前后比较。\par}
\clearpage
\pdfbookmark[1]{第 11 页：核心结果：全部确认略增，实质性均值近乎相同}{slide-11}
{\sffamily\bfseries\color{Navy}\fontsize{14}{19}\selectfont 第 11 页｜核心结果：全部确认略增，实质性均值近乎相同\par}
{\sffamily\color{Slate}\fontsize{8.5}{11}\selectfont Confirmed all rises; substantive means are nearly equal\par}
\IfFileExists{trace_q_slides.pdf}{\begin{center}\includegraphics[page=11,width=.76\linewidth]{trace_q_slides.pdf}\end{center}}{}
{\color{Slate}\fontsize{9}{12}\selectfont\textbf{讲解提示：}在 2.95→3.32 和 0.84→0.84 处稍停顿；指向森林图跨越零的区间。\par}
\vspace{1mm}
从最终工作簿逐条汇总，早期两年的已确认统计问题是 34 加 22，共 56 个；后期是 43 加 40，共 83 个。分别除以 19 篇和 25 篇论文，均值为 2.95 和 3.32，后期减早期约增加 0.37 个。这些是最终标签对应的实际计数。\par
这项差值的 95\% bootstrap 置信区间是负 1.32 到正 2.04，精确置换 P 值为 0.683。也就是说，点估计略高，但不确定性范围同时包含减少和增加。\par
实质性子集在早期是 11 加 5，共 16 个；后期是 11 加 10，共 21 个。两期论文数不同，换算后都约为每篇 0.84；差值约为负 0.002。对应区间为负 0.77 到正 0.74，精确置换 P 值为 1。\par
下面两条区间都跨过零，所以目前的数据没有建立明确的时期差异。同时，P 值很大，也不能证明两个时期完全相同，或者已经达到统计学上的等效。这里能直接说的是：在这个样本中，全部确认问题的平均数略高，而确认且实质的问题均值近乎相同。\par
{\color{Teal}\textbf{转场：}那么，全部确认口径里这约 0.37 个问题的增加，具体来自哪一部分？\par}
\clearpage
\pdfbookmark[1]{第 12 页：变化分解：呈现问题与受影响论文比例}{slide-12}
{\sffamily\bfseries\color{Navy}\fontsize{14}{19}\selectfont 第 12 页｜变化分解：呈现问题与受影响论文比例\par}
{\sffamily\color{Slate}\fontsize{8.5}{11}\selectfont Presentation issues account for the count difference\par}
\IfFileExists{trace_q_slides.pdf}{\begin{center}\includegraphics[page=12,width=.76\linewidth]{trace_q_slides.pdf}\end{center}}{}
{\color{Slate}\fontsize{9}{12}\selectfont\textbf{讲解提示：}先讲上表的加法关系，再切换到下表的“是否至少一个”；百分比差读成百分点。\par}
\vspace{1mm}
上半部分把全部确认问题拆成两类。确认的呈现问题从每篇 2.105 个增加到 2.480 个，差值约为正 0.375；确认且实质的问题从 0.842 变成 0.840，差值约为负 0.002。两部分相加，就是全部确认问题约正 0.373 的变化。因此，这个均值差在数值上来自呈现问题的增加，并不表示实质性问题也增加了。\par
下半部分换了一个问题：不是问每篇平均有多少个，而是问有多少篇至少出现一个。全部确认口径下，这个比例从 16 篇占 19 篇，也就是 84.2\%，变成 22 篇占 25 篇，也就是 88.0\%；增加了 3.8 个百分点。\par
确认且实质的口径下，比例从 42.1\% 变成 48.0\%，增加了 5.9 个百分点。但两项比例差的置信区间都很宽，P 值分别为 1 和 0.766，仍然没有建立明确的时期差异。\par
每篇平均问题数和至少有一个问题的论文比例，并不需要同步变化。一组问题可以分布在较少论文中，也可以分散在更多论文中。因此，实质性均值近乎相同，同时受影响论文比例略高，在数学上并不矛盾。\par
{\color{Teal}\textbf{转场：}最后还要解释一个口径差异：为什么论文的主要分析看起来是下降，而这里的确认且实质均值几乎不变。\par}
\clearpage
\pdfbookmark[1]{第 13 页：敏感性分析：纳入 probable 后为何出现下降}{slide-13}
{\sffamily\bfseries\color{Navy}\fontsize{14}{19}\selectfont 第 13 页｜敏感性分析：纳入 probable 后为何出现下降\par}
{\sffamily\color{Slate}\fontsize{8.5}{11}\selectfont The apparent decline comes from probable issues\par}
\IfFileExists{trace_q_slides.pdf}{\begin{center}\includegraphics[page=13,width=.76\linewidth]{trace_q_slides.pdf}\end{center}}{}
{\color{Slate}\fontsize{9}{12}\selectfont\textbf{讲解提示：}明确指出 4 个 probable 全在早期；把变化描述为口径敏感性，而非某种口径优劣。\par}
\vspace{1mm}
关键就在这一页的四个 probable，也就是很可能存在、但尚未确认的实质性问题。它们全部来自早期组。\par
如果只纳入确认且实质的问题，早期是 16 个，后期是 21 个，两期均值都约为 0.84。如果采用论文的主要结局，把很可能存在的实质性问题也算进去，早期就变成 16 加 4，合计 20 个，平均每篇 1.05 个；后期没有新增的 probable，所以仍然是 21 个，平均每篇 0.84 个。\par
这样，后期减早期的差值就从接近零，变成负 0.21。论文主要结果中这个表面上的下降，正是由纳入早期这四个 probable 条目产生的。它的 95\% 置信区间是负 1.02 到正 0.56，P 值是 0.654，仍然不足以建立时期差异。\par
这说明结论对证据确定程度的纳入标准有一定敏感性，两种口径都应透明报告。补充分析中，符合条件的实质性报告内部不一致是早期 1 个、后期 3 个；把它们和主要结局合并，均值为 1.11 对 0.96，P 值为 0.754，也没有改变整体解释。\par
{\color{Teal}\textbf{转场：}除了统计结果，我们也按各自的审查规则比较另外三个领域。\par}
\clearpage
\pdfbookmark[1]{第 14 页：其他三个领域：宽口径的前后比较}{slide-14}
{\sffamily\bfseries\color{Navy}\fontsize{14}{19}\selectfont 第 14 页｜其他三个领域：宽口径的前后比较\par}
{\sffamily\color{Slate}\fontsize{8.5}{11}\selectfont Extend the comparison to the other three review domains\par}
\IfFileExists{trace_q_slides.pdf}{\begin{center}\includegraphics[page=14,width=.76\linewidth]{trace_q_slides.pdf}\end{center}}{}
{\color{Slate}\fontsize{9}{12}\selectfont\textbf{讲解提示：}每格先报每篇均值，再说明方括号内是候选总数；三个领域不能都称为错误。\par}
\vspace{1mm}
其他三个领域也按最终工作簿 E 列的领域标签重新汇总。报告内部和计划材料的宽口径，要求 K 列是实质性不一致或轻微呈现不一致；保留但无法判断的条目不计入已建立问题。方案与报告领域的宽口径则统计全部保留差异，不能直接称为错误或违反方案。\par
报告内部已建立的问题从 3 条到 8 条，每篇均值从 0.16 到 0.32，差值约正 0.16，P 值 0.301。计划材料从 12 条到 17 条，每篇从 0.63 到 0.68，差值约正 0.05，P 值 0.893。方案与报告保留差异从 102 条到 112 条，但按论文数换算，均值从 5.37 降到 4.48，差值约负 0.89，P 值 0.256。\par
保留数与已建立数要分清：报告内部两期保留 4 条和 8 条，其中已建立 3 条和 8 条；计划材料保留 12 条和 20 条，其中已建立 12 条和 17 条。方案与报告保留差异中，解释字段有 2 条和 3 条为空；宽口径仍包括它们。三个均值差区间均跨过零。\par
这些新增结果按最终领域标签定义，也纳入原始计划相关关系，因此不同于论文原来只分析报告材料的口径。它们是事后描述性扩展；与统计结果一起看时，必须保留各领域不同的计数含义。\par
{\color{Teal}\textbf{转场：}下一页进一步只看各领域的实质性状态。\par}
\clearpage
\pdfbookmark[1]{第 15 页：其他三个领域：实质性结果与稀疏性}{slide-15}
{\sffamily\bfseries\color{Navy}\fontsize{14}{19}\selectfont 第 15 页｜其他三个领域：实质性结果与稀疏性\par}
{\sffamily\color{Slate}\fontsize{8.5}{11}\selectfont Substantive findings are sparse outside the statistical domain\par}
\IfFileExists{trace_q_slides.pdf}{\begin{center}\includegraphics[page=15,width=.76\linewidth]{trace_q_slides.pdf}\end{center}}{}
{\color{Slate}\fontsize{9}{12}\selectfont\textbf{讲解提示：}解释早期零并不证明真实发生率为零；区分新增领域口径与原报告材料口径。\par}
\vspace{1mm}
报告内部实质性问题从 1 条到 4 条，每篇均值从 0.05 到 0.16，P 值 0.370。计划材料实质性问题从 0 条到 4 条，每篇从零到 0.16，P 值 0.498；这四条集中在后期两篇论文。方案与报告的实质性差异从 83 条到 96 条，但每篇均值从 4.37 降到 3.84，P 值 0.446。\par
计划材料的早期观察值全为零，分层 bootstrap 无法从这些零里产生早期的非零值，所以区间下限为零。这并不证明真实发生率为零，也不能据此声称已经建立增长。样本稀疏时要同时看原始分布、重抽样局限和置换结果。\par
原报告材料口径中，实质性报告内部问题是早期 1 条、后期 3 条；按最终领域标签汇总则是 1 条和 4 条。多出的条目是 2024 年 jama.2024.11057 的 C003：最终工作簿将其归为报告内部不一致且具有实质性，来源对照表记录的原始关系是 P\_R。因此它进入最终领域比较，但不进入原报告材料口径。\par
方案与报告中的 substantive\_difference 仍是实质性差异的描述，不代表已经确认不当偏离。最终解释仍需要对应的计划版本、时间和理由。\par
{\color{Teal}\textbf{转场：}最后回到统计问题，看看总均值之下的机制组成有没有不同。\par}
\clearpage
\pdfbookmark[1]{第 16 页：统计机制的前后分布：总量相似不等于组成相同}{slide-16}
{\sffamily\bfseries\color{Navy}\fontsize{14}{19}\selectfont 第 16 页｜统计机制的前后分布：总量相似不等于组成相同\par}
{\sffamily\color{Slate}\fontsize{8.5}{11}\selectfont Which statistical mechanisms differ between publication cohorts?\par}
\IfFileExists{trace_q_slides.pdf}{\begin{center}\includegraphics[page=16,width=.76\linewidth]{trace_q_slides.pdf}\end{center}}{}
{\color{Slate}\fontsize{9}{12}\selectfont\textbf{讲解提示：}年份和实质性来自最终结果；机制来自另外的归档分类，不报告机制因果效应。\par}
\vspace{1mm}
这一页是补充机制分类的时期分布，不是直接从最终工作簿读取的新增字段。候选的纳入、年份和实质性采用最终结果；主要机制来自先前归档的 Notes 分类。每格先报候选数，再报每篇均值，早期分母为 19，后期为 25。\par
全部确认中，人数与分母从 13 条到 26 条，每篇从 0.68 到 1.04；分析或测量定义从 3 条到 17 条，每篇从 0.16 到 0.68。区间与分布摘要则从 14 条到 10 条，每篇从 0.74 降到 0.40。说明两个时期的机制组成不同，即使总均值变化不大，也可能有方向相反的分项变化。\par
实质性部分同样有抵消：分析定义从零到 3 条，效应量从 3 条到 5 条，而区间从 4 条到 3 条，推断从 3 条到 1 条。合计仍是早期 16 条、后期 21 条，换算后都约为每篇 0.84。\par
这些机制是事后分类，部分记录还有分类歧义；本页没有另做类别层面的显著性检验。可以描述归档分类中的组成变化，但不能称为最终审查工作簿已经独立验证的机制差异，更不能归因于疫情。只有 139、37 及其按年份和论文分组的审查结果直接以最终标签为准。\par
{\color{Teal}\textbf{转场：}最后用四点收束这次汇报。\par}
\clearpage
\pdfbookmark[1]{第 17 页：结束：把证据、判断、机制与时期比较连接起来}{slide-17}
{\sffamily\bfseries\color{Navy}\fontsize{14}{19}\selectfont 第 17 页｜结束：把证据、判断、机制与时期比较连接起来\par}
{\sffamily\color{Slate}\fontsize{8.5}{11}\selectfont Keep evidence, mechanism and period interpretation distinct\par}
\IfFileExists{trace_q_slides.pdf}{\begin{center}\includegraphics[page=17,width=.76\linewidth]{trace_q_slides.pdf}\end{center}}{}
{\color{Slate}\fontsize{9}{12}\selectfont\textbf{讲解提示：}强调两个核心均值、另外三个领域、可追溯语料和解释边界。\par}
\vspace{1mm}
第一，全部确认的统计问题均值从每篇 2.95 到 3.32，约增加 0.37，数值变化来自呈现问题。确认且实质的均值则约为 0.84 对 0.84。区间仍较宽，不能说已证明没有变化或达到等效。\par
第二，另外三个领域已经分别补充前后比较，但都没有建立明确的时期差异。报告不一致、计划材料问题，以及方案报告差异，各自的定义必须保留，尤其不能把所有方案差异都称为错误。\par
第三，本次以四份 completed 工作簿为结果依据：440 条候选中保留 413 条、排除 27 条，全部完成。codebook 把处置、确定性和实质性分开，只有先说明筛选条件，139、37 和 41 这几个数量才可解释。\par
第四，最终工作簿中保留条目的 Notes 为空。七类机制和两个详细例子来自另外归档的原始说明，是补充分析；本次核对了它们与最终候选及实质性标签的一致性。这些结果属于探索性样本分析，不能估计所有真实问题的漏检率，也不能说明疫情造成了变化。谢谢大家。\par
\end{document}
